% Curriculum Vitae — Steven Hoekstra
% Build:  pdflatex -jobname=Steven_Hoekstra_CV cv.tex   (run from the cv/ directory)
\documentclass[11pt,a4paper]{article}

\usepackage[T1]{fontenc}
\usepackage[utf8]{inputenc}
\usepackage[left=2.05cm,right=2.05cm,top=2.05cm,bottom=2.05cm]{geometry}
\usepackage{mathpazo}
\usepackage[scaled=0.92]{helvet}
\usepackage{microtype}
\usepackage{xcolor}
\usepackage{tabularx}
\usepackage{needspace}
\usepackage{fancyhdr}
\usepackage[hidelinks]{hyperref}

\definecolor{rule}{HTML}{D0D6DC}
\definecolor{accent}{HTML}{243F56}
\definecolor{muted}{HTML}{48535D}
\hypersetup{
  pdftitle={Curriculum Vitae - Steven Hoekstra},
  pdfauthor={Steven Hoekstra},
  colorlinks=true, urlcolor=accent, linkcolor=accent
}

\setlength{\parindent}{0pt}
\setlength{\parskip}{0pt}
\setlength{\emergencystretch}{1em}
\pagestyle{fancy}
\fancyhf{}
\fancyfoot[L]{\sffamily\scriptsize\color{muted}Steven Hoekstra\enspace\textperiodcentered\enspace Curriculum Vitae}
\fancyfoot[C]{\sffamily\scriptsize\color{muted}13 October 2026}
\fancyfoot[R]{\sffamily\scriptsize\color{muted}\thepage}
\renewcommand{\headrulewidth}{0pt}
\renewcommand{\footrulewidth}{0pt}

% Dates have a fixed column; titles wrap independently and align at the top.
\newcommand{\cvsection}[1]{%
  \par\Needspace{4\baselineskip}\vspace{0.8em}%
  {\sffamily\bfseries\color{accent}\fontsize{11.8}{14}\selectfont #1\par}%
  \vspace{0.2em}{\color{rule}\hrule height 0.35pt}\vspace{0.4em}%
}
\newcommand{\entry}[2]{%
  \par\Needspace{3\baselineskip}\vspace{0.35em}%
  \noindent\begin{minipage}[t]{\dimexpr\linewidth-3cm-0.8em\relax}%
    \vspace{0pt}\raggedright\fontsize{10.6}{12.3}\selectfont\textbf{#1}\par
  \end{minipage}\hfill
  \begin{minipage}[t]{3cm}%
    \vspace{0pt}\raggedleft\fontsize{9.5}{11.2}\selectfont\textcolor{muted}{#2}\par
  \end{minipage}\par
}
\newcommand{\titleonly}[1]{%
  \par\Needspace{3\baselineskip}\vspace{0.35em}%
  {\fontsize{10.6}{12.3}\selectfont\bfseries\raggedright #1\par}%
}
\newcommand{\note}[1]{\par{\fontsize{9.8}{11.5}\selectfont\color{muted}\raggedright #1\par}}
\newcommand{\subhead}[1]{\par\Needspace{3\baselineskip}\vspace{0.5em}{\sffamily\small\bfseries\color{muted}#1\par}\vspace{0.1em}}
\newenvironment{cvitems}{%
  \par\begingroup\fontsize{9.8}{11.5}\selectfont\color{muted}%
  \begin{list}{\textbullet}{%
    \setlength{\leftmargin}{1.2em}\setlength{\labelwidth}{0.5em}%
    \setlength{\labelsep}{0.7em}\setlength{\topsep}{0.3em}%
    \setlength{\partopsep}{0pt}\setlength{\itemsep}{0.12em}\setlength{\parsep}{0pt}}%
  \raggedright
}{\end{list}\endgroup}

\begin{document}

% --- header ----------------------------------------------------------------
{\fontsize{30}{34}\selectfont\sffamily\bfseries\color{accent}Steven Hoekstra}\\[0.45em]
{\small
PhD Candidate in Economics, University of Amsterdam \& Tinbergen Institute\\
\href{mailto:s.g.hoekstra@uva.nl}{s.g.hoekstra@uva.nl} \quad$\cdot$\quad
\href{https://sghoekstra.github.io}{sghoekstra.github.io} \quad$\cdot$\quad
\href{https://github.com/SGHoekstra}{github.com/SGHoekstra}
}
\par\vspace{0.3em}

\cvsection{Profile}
I study how heterogeneous households and firms, boundedly rational expectations and market
interactions shape monetary and fiscal policy transmission. I connect heterogeneous-agent
equilibrium models and agent-based models, using neural-network solvers, simulation-based
inference and microdata to assess economic mechanisms and the reliability of policy predictions.

\cvsection{Research Interests}
Macroeconomics, heterogeneous-agent models, agent-based modelling, expectations and bounded
rationality, policy evaluation, networks, simulation-based inference, high-performance computing.

\cvsection{Education}
\entry{PhD in Economics \textnormal{--- University of Amsterdam}}{2023 --- current}
\note{Supervisors: Prof.\ dr.\ C.H. Hommes and Prof.\ dr.\ F.P. Pijpers}
\note{Expected thesis submission: May 2027; expected defence: September 2027}
\entry{MPhil Economics and Econometrics \textnormal{--- Tinbergen Institute} \textit{(cum laude)}}{2021 --- 2023}
\note{Specialisation: Complexity, Mathematics, and Computation}
\note{Thesis: An Agent Based Model of the Dutch Economy}
\note{Supervisors: Prof.\ dr.\ C.H. Hommes and Prof.\ dr.\ F.P. Pijpers}
\entry{MSc Econometrics \textnormal{--- University of Amsterdam} \textit{(cum laude)}}{2020 --- 2021}
\note{Specialisation: Complexity and Economic Behaviour}
\note{Thesis: Endogenous Financial Instability via Heterogeneous Beliefs}
\note{Supervisor: Prof.\ dr.\ C.H. Hommes}
\entry{MSc Economics \textnormal{--- London School of Economics}}{2019 --- 2020}
\note{Specialisation: Public Economics}
\note{Thesis: The Inequality of Education using a Two Agent New Keynesian Model}
\note{Supervisor: Prof.\ dr.\ F. Cowell}
\entry{Exchange Programme \textnormal{--- Royal Melbourne Institute of Technology}}{2018 --- 2019}
\note{Minor in Finance and Programming}
\entry{BSc Economics and Business Economics \textnormal{--- University of Groningen}}{2016 --- 2019}
\note{Thesis: The Impact of News Coverage on Security Prices}
\entry{BSc Econometrics \textnormal{--- University of Groningen}}{2015 --- 2018}
\note{Thesis: Arbitrage Smoothing and Term Structure Models}

\cvsection{Working Papers and Research Software}
\titleonly{Forecasting with a Dutch Agent-Based Model and Optimally Behaving Households}
\note{Working paper, with Cars Hommes and Frank Pijpers. \href{https://sghoekstra.github.io/papers/forecasting-dutch-abm.pdf}{PDF}}
\note{Runner-up, Society for Computational Economics Graduate Student Paper Award, CEF 2026}
\titleonly{BeforeIT.jl: High-Performance Agent-Based Macroeconomics Made Easy}
\note{Research software; package co-author. Developed the validation suite.}
\note{\href{https://github.com/bancaditalia/BeforeIT.jl}{Package} $\cdot$ \href{https://arxiv.org/abs/2502.13267}{Original preprint (2025)}}

\newpage
\cvsection{Work in Progress}
\titleonly{When Do Estimated Macroeconomic Models Misread Wage--Price Transmission?}
\note{Single-author. Known-truth experiments on misspecification and intervention-response reliability;
Dutch wage--price transmission and specification sensitivity during the 2022 inflation episode.}
\titleonly{Shortage or Misallocation? The Welfare Loss from Dutch Housing Policy, 2015--2035}
\note{Integrating housing, labour mobility, demography and firm dynamics in a macroeconomic
ABM; calibration and welfare comparisons are ongoing.}

\cvsection{Work Experience}
\entry{De Nederlandsche Bank --- PhD Intern}{Apr.\ --- Dec.\ 2025}
\begin{cvitems}
\item Implemented the Bank of Canada CANVAS agent-based model for the Netherlands
\item Estimated inflation channels using neural ratio estimation
\item Wrote a note on the wage-price channel using this estimated model
\end{cvitems}
\entry{Statistics Netherlands (CBS) --- Thesis Intern}{Mar.\ --- Sept.\ 2023}
\begin{cvitems}
\item Calibrated the agent-based model of Poledna et al.\ (2023) with Dutch register and national accounts data
\item Developed an extension incorporating optimal consumption
\item Ran simulations with both models and analysed their internal dynamics
\end{cvitems}
\entry{SEO Amsterdam Economics --- Research Assistant}{2020 --- 2021}
\begin{cvitems}
\item Researched academic pressure among secondary school students
\item Conducted stakeholder interviews and set up a nationwide survey
\item Performed statistical analysis on the survey outcomes
\end{cvitems}

\cvsection{Grants}
\titleonly{NWO Compute Grants --- Snellius}
\note{Three awarded compute grants of 1,000,000 CPU hours each; 3,000,000 CPU hours in total.}

\cvsection{Teaching}
\subhead{Graduate}
\entry{Behavioural Macro and Finance --- Lecturer}{2025}
\note{Developed and taught the MSc lecture series with Cars Hommes, University of Amsterdam.}
\titleonly{MPhil guest lectures --- Tinbergen Institute}
\note{Annual guest lectures in Behavioural Macroeconomics and Complexity.}
\subhead{Undergraduate}
\titleonly{BSc Econometrics thesis supervision --- University of Amsterdam}
\note{Ongoing.}
\entry{Mathematics 3: Advanced Linear Algebra --- Tutor}{2025 --- 2026}
\entry{Mathematics 2: Linear Algebra --- Tutor}{2024 --- 2025}
\entry{Microeconomics for AE --- Tutor}{2023 --- 2024}

\newpage
\cvsection{Presentations}
\begin{tabularx}{\linewidth}{@{}p{1.6cm}X@{}}
2026 & Computing in Economics and Finance (CEF), Society for Computational Economics\\[0.15em]
2026 & Workshop on Economic Science with Heterogeneous Interacting Agents (WEHIA)\\[0.15em]
2025 & Agent-Based Modelling for Policy (ABM4Policy), European Central Bank, Frankfurt\\[0.15em]
2025 & 7th Behavioral Macroeconomics Workshop, University of Bamberg\\[0.15em]
2025 & MacroABM 2nd Workshop, \textit{Microdata in Macromodels}, Vienna University of Economics and Business\\[0.15em]
2025 & Workshop on Economic Science with Heterogeneous Interacting Agents (WEHIA)\\[0.15em]
2025 & De Nederlandsche Bank, Amsterdam\\[0.15em]
2023 & De Nederlandsche Bank, Amsterdam\\
\end{tabularx}

\cvsection{Honours and Scholarships}
\subhead{Honours}
\entry{Cum Laude --- MPhil Economics and Econometrics, Tinbergen Institute}{2023}
\entry{Cum Laude --- MSc Econometrics, University of Amsterdam}{2021}
\subhead{Scholarships}
\entry{Full scholarship (merit based) --- Tinbergen Institute}{2022}
\entry{Tuition fee waiver --- Tinbergen Institute}{2021 --- 2023}

\cvsection{Visits and Summer Schools}
\entry{Visiting Scholar --- IIASA, Laxenburg}{Aug.\ 2023}
\note{International Institute for Applied Systems Analysis. Visited Dr.\ S. Poledna to work on agent-based models}
\entry{Tools for Macroeconomists: Advanced Tools --- University of Oxford}{Aug.\ 2021}
\note{Doctoral-level course on DSGE models: solving and simulating models with heterogeneous
agents, continuous time models, occasionally binding constraints and bounded rationality}
\entry{Tools for Macroeconomists: The Essentials --- University of Oxford}{Aug.\ 2021}
\note{Doctoral-level course on DSGE models: numerical analysis, perturbation and projection
methods, Blanchard-Kahn conditions and sunspot solutions, the Kalman filter and Bayesian
estimation}

\cvsection{Skills}
\begin{tabularx}{\linewidth}{@{}>{\raggedright\arraybackslash}p{5.1cm}>{\raggedright\arraybackslash}X@{}}
Programming & Matlab, Python, R, Julia, C, SQL, Bash\\[0.15em]
High-performance computing & Parallel computing (MPI, BSPlib, Slurm)\\[0.15em]
Methods & Agent-based modelling, heterogeneous-agent models, neural-network solvers,
simulation-based inference, network methods, machine learning, econometrics,
dynamic stochastic general equilibrium modelling (Dynare)\\
\end{tabularx}

\cvsection{References}
\begin{tabularx}{\linewidth}{@{}XX@{}}
\textbf{Prof.\ dr.\ C.H. Hommes} & \textbf{Prof.\ dr.\ F.P. Pijpers}\\
{\small CeNDEF, Faculty of Economics,} & {\small Korteweg-de Vries Institute,}\\
{\small University of Amsterdam} & {\small Faculty of Science, University of Amsterdam}\\
{\small Bank of Canada, Tinbergen Institute} & {\small Statistics Netherlands}\\
{\small \href{mailto:c.h.hommes@uva.nl}{c.h.hommes@uva.nl}} &
{\small \href{mailto:f.p.pijpers@uva.nl}{f.p.pijpers@uva.nl}}\\
\end{tabularx}

\end{document}
